\section{Audit del benchmark Stage D e definizioni delle metriche}
\label{app:auditstaged}

\subsection{Tracciabilità dell'inferenza}

Il benchmark ha impiegato i checkpoint Stage A e Stage D identificati dai rispettivi hash SHA-256 registrati negli snapshot degli esperimenti. L'audit del call graph ha verificato che l'adattatore di feature non riceve il problema manufatto esatto: la GNN opera sulle coordinate, sul campo FEM corrente, sulle sue derivate discrete e sull'identificatore di bordo. I target di riferimento non vengono caricati durante l'inferenza; non vengono passati alla rete parametri fisici espliciti o stati futuri. La conclusione di assenza di leakage riguarda il percorso di inferenza e i canali effettivamente verificati.

\subsection{Filtro geometrico: dati Stage D e limitazioni}

Nel benchmark Stage D il filtro interviene in 211 finestre su 280 (75.36\%) e $\alpha$ medio è 0.4968. Per Stage A gli interventi sono 265 su 280 (94.64\%) e $\alpha$ medio è circa 0.175. Il minimo post-filtro registrato per Stage D è $9.999999999999865\times10^{-5}$; il numero di celle invertite e degeneri dopo il filtro è zero. La minima differenza stretta rispetto a $\epsilon=10^{-4}$ è dell'ordine di $10^{-18}$. Alcune colonne del vecchio CSV temporale relative alle metriche raw e post-filtro sono duplicate; il numero di celle raw non ammissibili non è pertanto ricostruibile per tutti gli eventi Stage D.

\subsection{D-GCL e discrepanza degli incrementi di massa}

I massimi residui D-GCL del benchmark Stage D sono:
\[
R_{\mathrm{GCL,alg}}^\infty=1.9665\times10^{-16},\qquad
R_{\mathrm{GCL,elem}}^{\max}=4.6972\times10^{-17}.
\]
Nel log storico dei test, nove test risultano superati e uno non supera la soglia stretta del free-stream: l'errore $L^\infty$ è $2.517\times10^{-12}$ rispetto alla tolleranza $10^{-12}$. Questo esito non coincide con i residui algebrico ed element-wise, che sono diagnostiche differenti. La verifica indipendente E1, descritta nell'Appendice~\ref{app:stageef}, misura inoltre il residuo del solve e la tolleranza GMRES effettiva; non viene usata per riscrivere retroattivamente il log Stage D.

La discrepanza degli incrementi di massa è
\[
D_M^n=\abs{(M_h^n-M_h^{n-1})-(M_{\mathrm{ref}}^n-M_{\mathrm{ref}}^{n-1})},
\]
con versione relativa normalizzata tramite $\max(\abs{M_{\mathrm{ref}}^0},10^{-12})$. La media relativa è circa $1.62\times10^{-4}$ per Fixed FEM e $6.36\times10^{-3}$ per Stage D. La quantità non rappresenta il residuo completo dell'identità integrale della PDE; non include esplicitamente in un'unica identità tutti i contributi di forzante e flussi al contorno. Il dato di Dirichlet è $g_D\equiv0$; la massa varia per la forzante e per il flusso diffusivo normale. Per tale ragione la quantità viene denominata discrepanza degli incrementi di massa e non errore di conservazione.
